\section{Project Overview and Current Evidence Baseline}
\label{sec:project_overview_current_baseline}

This dissertation is titled \emph{Schema Validity Is Not Enough: Zero-Trust Evaluation of Local SLMs for Industrial Robot Task Planning with Microsoft Foundry Local}. The project investigates whether local small language models (SLMs), served through Microsoft Foundry Local, can support industrial robot task-planning workflows when their outputs are treated as untrusted action proposals rather than executable commands. The central claim is deliberately bounded: within the evaluated benchmark and under the defined validation policy, local SLMs can contribute to robot task planning only when deterministic validation mediates every model-generated proposal before any downstream execution context is reached.

The work is positioned in industrial robotics and manufacturing automation. In the target use case, a factory operator gives a natural-language instruction to a local workcell assistant. A local SLM proposes a structured robot action, and a zero-trust validation layer evaluates whether the proposal is parseable, JSON-valid, schema-valid, semantically valid, safety-valid and execution-eligible. This separation is essential because a schema-valid output can still be semantically wrong, unsafe, ambiguous or unsupported by the current workcell state.

Microsoft Foundry Local is used as the local inference framing for the project. It provides a practical route for evaluating local model serving, privacy and offline-resilience trade-offs without treating the model itself as a trusted controller \cite{microsoft_foundry_local_docs}. The dissertation therefore evaluates Foundry Local as part of a local-first evidence chain, not as a production-certified robot-control platform.

Prototype 5 is the final evidence orchestration and reporting layer. It consolidates prior evidence, records claim boundaries, generates the final claims matrix and manifest, and adds Prototype 5 evidence for custom quantisation, Phi-family local responses, local-vs-cloud comparison and live local resource profiling. Prototype 1 is retained only as optional early feasibility context; its absence from core audit evidence is a contextual documentation gap, not a missing proof for the final claims.

The current evidence baseline is complete for the dissertation pack: the Prototype 5 test suite passed, the final orchestrator completed, and the final dissertation pack generator completed. The evidence remains benchmark-based and prototype-scoped. It does not establish real-world robot safety, factory deployment readiness, clinical or production-grade assurance, or statistical generalisation beyond the evaluated command set.
